\chapter{Quantitative inference with certified nonlinear bounds}
\label{chap:ci-nonlinear}

A positive contrast answers whether the specified equilibrium comparison class
can explain the observations. Its quadratic entropy bound can discard much of
the available probability information. This chapter gives a nonlinear replacement
and an independently checkable bracket on the resulting optimization problem.
It also identifies a second source of looseness: even an exact observational
KL projection need not determine physical dissipation. An additional calibration
of microscopic jump activity can sharpen that separate physical step.

The variational representation of relative entropy, convex information
projection and likelihood-score testing are established ingredients. The
modern composite-null testing literature gives closely related log-optimal
score constructions\cite{ciLarsson2025,ciGrunwald2024,ciLarsson2026}. The
contribution here is their explicitly certified use with a calibrated continuous
moment class, a nonrectangular finite-data confidence region, and a separate
physical-resource calculation. Simulations and numerical optimization do not
establish the physical premises.

\section{The comparison problem and its scope}
Let $k_1,\ldots,k_m$ be positive integer lag indices, $w_i>0$ sum to one, and
$t_i>0$ bound the full expected path duration associated with each target
record. Write $p_i$ for the probability that the two binary readings agree.
Suppose a valid joint confidence event places the target vector in
\[
 {\cal P}=\{p\in\textstyle\prod_i[\ell_i,h_i]:a\cdot p\ge b\}.
\]
The contrast $a$ must have been declared before its test records, or selected
using independent pilot data. The halfspace can use the one-sided conditional
Hoeffding bound proved in the preceding work; the default box is $[0,1]^m$.
An empty region is rejected as an inconsistent input, not interpreted as a
scientific certificate.

Let $V$ be rational vertices of a calibrated polygon in
$\{(u,A):0\le u\le A\le1\}$. Here $u$ bounds the square of the stationary
detector mean and $A$ bounds the stationary second moment of its conditional
mean. Define the closed reversible comparison class
\[
 {\cal K}_V=\operatorname{conv}\left\{
  q_i={1+u+(A-u)x^{k_i}\over2}:(u,A)\in V,\ 0\le x\le1\right\}.
\]
The preceding calibration chapter explains the access needed to measure $A$.
The convex-hull representation allows distinct spectral atoms and arbitrary
calibration points inside the polygon. Indeed, a convex combination has
effective $u=\sum_j\theta_j u_j$, $A=\sum_j\theta_j A_j$, and spectral weights
proportional to $\theta_j(A_j-u_j)$. Conversely, decomposing $(u,A)$ into
vertices gives such a convex mixture. The atom $x=0$ is part of the closed
comparison class and may represent an infinite-rate limit; it need not be
an exact finite-rate physical realization.

The quantity to bracket is
\begin{equation}
 B({\cal P},{\cal K}_V;w)=
 \inf_{p\in{\cal P},q\in{\cal K}_V}
       \sum_i w_i\,\operatorname{kl}(p_i\Vert q_i),
 \label{eq:ci-nonlinear-projection}
\end{equation}
where $\operatorname{kl}$ is Bernoulli relative entropy in nats. For a
stationary ordinary-reversal CTMC, the earlier additive-reversibilization
bridge gives
\[
 \sigma\ge {4B\over\sum_iw_it_i}.
\]
This comparison class generally contains reference models that cannot be the
additive reversibilization of the same hidden generator producing $p$.
Consequently, $B$ is an observational relaxation. A feasible upper bound on
$B$ is not an upper bound on physical entropy production, or proof that its
lower bound is physically attainable. Odd-variable physical reversal,
measurement disturbance and an uncalibrated initial transient remain outside
this stationary interpretation.

\section{An exact dual certificate and a primal bracket}
Choose arbitrary positive rational scores $h_{i0},h_{i1}$. Let
\begin{align*}
 L(h)&=\inf_{p\in{\cal P}}\sum_iw_i
       [(1-p_i)\log h_{i0}+p_i\log h_{i1}],\\
 U(h)&=\sup_{q\in{\cal K}_V}\sum_iw_i
       [(1-q_i)h_{i0}+q_i h_{i1}].
\end{align*}
Then
\begin{equation}
 B\ge\max\{0,L(h)-\log U(h)\}.
 \label{eq:ci-nonlinear-dual}
\end{equation}
For each fixed pair $(p,q)$, regard $(i,y)$ as a joint record with probabilities
$w_i\operatorname{Ber}(p_i)_y$ and $w_i\operatorname{Ber}(q_i)_y$.
The relative-entropy variational inequality gives
$D(P\Vert Q)\ge E_P\log h-\log E_Qh$.
Taking the worst cases over the two convex sets proves
\eqref{eq:ci-nonlinear-dual}.

Both optimization steps in a supplied certificate are checked on their whole
domains. For a calibration vertex $(u,A)$, the normalization is the maximum on
$[0,1]$ of the rational polynomial
\[
 \sum_iw_i\left[h_{i0}+(h_{i1}-h_{i0})
       {1+u+(A-u)x^{k_i}\over2}\right].
\]
Rational Bernstein coefficients after interval subdivision give a rigorous
upper bound. No sampling grid is used as a proof of this maximum.
The logarithms are enclosed by the inherited rational atanh-series routine.
Substituting their lower endpoints gives a linear lower objective on
${\cal P}$. Its minimum is an exact fractional-knapsack calculation: begin
at a box minimizer, and if the contrast constraint fails, move coordinates in
the direction increasing $a\cdot p$, ordered by objective cost per unit
increase. These rational moves attain the minimum of the enclosed linear
objective. The continuous normalization and outward logarithm enclosures then
give a rational lower bound on $B$.

Conversely, supply rational $p\in{\cal P}$ and finitely many rational atoms
$(v_j,x_j,\theta_j)$ with $v_j\in V$, $x_j\in[0,1]$, $\theta_j\ge0$ and
$\sum_j\theta_j=1$. They produce an exactly feasible $q\in{\cal K}_V$.
Upper logarithm enclosures give
\[
 B\le\sum_iw_i\operatorname{kl}(p_i\Vert q_i)\le B_{\rm upper}.
\]
Thus the reported gap $B_{\rm upper}-B_{\rm lower}$ measures the remaining
uncertainty of the relaxation calculation. The numerical optimizer proposes
scores and mixtures; the rational checks retain their validity even when it
fails to converge or uses a coarse proposal grid.

For an interior minimizing pair $(p^*,q^*)$, the dual has no gap. Set
$h_{i1}=p_i^*/q_i^*$ and
$h_{i0}=(1-p_i^*)/(1-q_i^*)$. First-order optimality in $q$ implies
$E_Qh\le1$ for all comparison vectors, with equality at $q^*$.
First-order optimality in $p$ implies
$E_P\log h\ge E_{P^*}\log h$ throughout ${\cal P}$. Substitution gives
equality in \eqref{eq:ci-nonlinear-dual}. Boundary minimizers can require
limiting scores; the implementation claims only the brackets it actually
checks, not attainment at finite scores in every case.

For a fixed alternative, normalizing $h$ by $U(h)$ gives a nonnegative score
with expectation at most one throughout the composite equilibrium null.
Its expected logarithm is the familiar information-projection testing
objective. This connection does not make a score fitted to its own data an
automatically valid significance test. Here data-dependent scores are safe
because they are minimized over a previously valid confidence region that
contains the entire true probability vector on the same event. Calibration
failure probabilities must be included in that event separately.

\section{A realizable calibration example}
Consider a stationary twenty-state ring with conditional detector mean
$f_j=\cos(2\pi j/20)$. Its uniform stationary law gives $m=0$, $A=1/2$.
Put $r=\sqrt{497}/25$, $\omega=\arccos(19/\sqrt{497})$,
$\gamma=-\log r$, and $\phi=2\pi/20$. Choose clockwise and anticlockwise
rates
\[
 a={1\over2}\left({\gamma\over1-\cos\phi}
                   +{\omega\over\sin\phi}\right),\qquad
 b={1\over2}\left({\gamma\over1-\cos\phi}
                   -{\omega\over\sin\phi}\right).
\]
Both are positive: numerically $a\simeq2.06125923$ and $b\simeq0.27982488$.
Independent binary readings with means $f_j$ have agreement probability
$p(t)=1/2+e^{-\gamma t}\cos(\omega t)/4$. The specified constants give
$p(1)=0.69$, $p(2)=0.59$ exactly. The predeclared contrast
$(38/25,-1)$ therefore takes the value $1147/2500=0.4588$.
Its calibrated reversible support is
$19/25-1/2+(19/25)^2/4=0.4044$; allowing $A$ to rise to one raises
that support to $0.5488$. Positive and decreasing correlations alone do not
expose the discrepancy; their scale relative to calibrated detector variance
does.

Using the idealized confidence halfspace with threshold $0.4588$, equal
mixture weights and durations $(1,2)$, the saved deterministic calculation
brackets the calibrated record-KL relaxation between
$0.0010251000031$ and $0.0010265638733$. The uncalibrated bracket is
$[0,6.43\times10^{-14}]$. These are model calculations, not finite-record
experimental confidence statements. The true model entropy rate is about
$3.55736$, whereas the certified lower rate is about $0.002734$.
The narrow optimization bracket therefore reveals a substantial evidential
and physical-bridge limitation, not an unresolved numerical calculation.

\section{The physical bridge has its own limit}
For an unordered edge, put
$a=\pi_iq_{ij}$, $b=\pi_jq_{ji}$, $c=a+b$, and
$x=|a-b|/(a+b)$. Edges with zero traffic contribute zero. For finite entropy
assume positive reverse flux whenever forward flux is positive. Define
\[
 j(x)={ (1+x)\log(1+x)+(1-x)\log(1-x)\over2},
 \qquad e(x)=x\log{1+x\over1-x}.
\]
The path-KL rate from $Q$ to $S=(Q+Q^*)/2$ and the entropy-production rate are
exactly
\[
 d=\sum_{\{i,j\}}c_{ij}j(x_{ij}),\qquad
 \sigma=\sum_{\{i,j\}}c_{ij}e(x_{ij}).
\]
The escape rates of $Q$ and $S$ agree pointwise, so there is no omitted
waiting-time term in $d$. Total stationary microscopic jump activity is
$K_0=\sum c_{ij}$.

As $x\to1$, $j(x)\to\log2$ while $e(x)\to\infty$. A fixed-activity driven
ring therefore has bounded path-KL rate but unbounded entropy production.
An arbitrarily accurate endpoint projection cannot eliminate this physical
loss of information. The data do not determine a finite upper bound on total
dissipation without further assumptions.

There is, however, a sharper lower bridge if a separate physical calibration
gives $K_0\le K$. Write $G=e\circ j^{-1}$ on $[0,\log2)$. Then
\begin{equation}
 \sigma\ge K\,G\left({B\over K\bar t}\right),
 \qquad \bar t=\sum_iw_it_i,
 \quad 0\le {B\over K\bar t}<\log2.
 \label{eq:ci-activity-bridge}
\end{equation}
To prove it, $j'(x)=\operatorname{atanh}x$ and
\[
 {e'(x)\over j'(x)}
 =2+{2x\over(1-x^2)\operatorname{atanh}x}
\]
increase strictly for $0<x<1$. The derivative of the second term has the
sign of $(1+x^2)\operatorname{atanh}x-x$, which is positive because its
value at zero is zero and its derivative is
$2x\operatorname{atanh}x+2x^2/(1-x^2)>0$.
Thus $G$ is increasing and convex, with $G(0)=0$ and $G'(0)=4$.
Jensen's inequality gives $\sigma\ge K_0G(d/K_0)$.
The perspective $K\mapsto KG(d/K)$ is nonincreasing: convexity gives
$G(z)-zG'(z)\le0$. Finally data processing gives $d\ge B/\bar t$.
These three inequalities prove \eqref{eq:ci-activity-bridge}.
It improves $\sigma\ge4B/\bar t$, with strict improvement for nonzero
evidence and finite activity cap except for numerical rounding.

The activity-calibrated envelope is sharp for full-path observations of a
uniform ring with clockwise and anticlockwise rates
$K(1+x)/2$ and $K(1-x)/2$: every edge has the same bias, the stationary law is
uniform, total activity equals $K$, and Jensen is an equality. This is a
sharpness statement about the full-path physical bridge. Restricted endpoint
observations usually lose information, and the moment relaxation can lose
more; no endpoint sharpness follows from the ring example.

After normalization by activity, $d/K_0$ is the Jensen--Shannon divergence
between the stationary directed-flux distribution and its transpose, while
$\sigma/K_0$ is half their Jeffreys divergence when Jeffreys is defined as
the sum of the two directed KL divergences. Bounds between these divergences
and general finite-support optimization methods are established
\cite{ciCrooks2008,ciFalasco2022,ciGuntuboyina2014}. The explicit envelope
above is proved here for this inference contract; historical priority of the
scalar divergence inequality is not claimed.

The implementation inverts $j$ by rational bisection. If an upper logarithm
enclosure proves $j(x_-)\le B/(K\bar t)$, then
$Ke(x_-)$ with a lower logarithm enclosure is a certified lower entropy
bound. The previous linear bound is retained if it is numerically larger.
If the supplied evidence exceeds $K\bar t\log2$, the routine reports an
incompatible joint contract rather than an entropy estimate. The microscopic
activity cap cannot be inferred from a count of coarse detector changes:
unobserved jumps still contribute activity. Its measurement cost, preparation
cost and timing calibration belong in the experimental resource account.

\section{A scoped transfer to overdamped diffusions}
The observational certificate is not restricted to a finite state space.
Consider the stationary diffusion
\[
 dX_t=b(X_t)\,dt+\sqrt{2D}\,dW_t
\]
on a connected compact periodic domain, with smooth drift, constant symmetric
positive-definite $D$, and smooth strictly positive stationary density $\pi$.
All state coordinates have even physical time-reversal parity. Define
\[
 b_s=D\nabla\log\pi,\qquad v=b-b_s,\qquad
 \sigma=\int v^TD^{-1}v\,\pi\,dx.
\]
This is the standard steady current expression for overdamped entropy
production under the stated physical convention\cite{ciSeifert2012}.
The reversible comparison diffusion has drift $b_s$ and the same diffusion
matrix and initial law $\pi$.

The drift likelihood formula\cite{ciGirsanov1960} gives
\begin{equation}
 D(\mathbb P_b^{[0,t]}\Vert\mathbb P_{b_s}^{[0,t]})
   ={1\over4}E_b\int_0^t v(X_s)^TD^{-1}v(X_s)\,ds
   ={t\sigma\over4}.
 \label{eq:ci-diffusion-bridge}
\end{equation}
The factor follows from the noise covariance $2D$: the control in the
Brownian likelihood is $(\sqrt{2D})^{-1}v$, and relative entropy is one half
its expected squared norm. On the compact periodic domain the drift
difference is bounded, so Novikov's condition and finite energy hold.
On a noncompact domain the conclusion instead requires nonexplosion,
the appropriate boundary or decay conditions, a valid true likelihood
martingale and finite expected energy. Novikov is a sufficient condition,
not a claim automatically valid for every linear example and every horizon.

The reference generator is
$L_s=\pi^{-1}\nabla\!\cdot(\pi D\nabla)$.
Periodic integration by parts makes it self-adjoint and nonpositive in
$L^2(\pi)$. For a fixed bounded detector conditional mean $f$, the spectral
theorem therefore gives
\[
 E_s[f(X_0)f(X_t)]
   =m^2+\int_{[0,\infty)} e^{-\lambda t}\,d\nu(\lambda),
 \quad \nu\ge0,\quad\nu([0,\infty))=A-m^2,
\]
where $m=E_\pi f$ and $A=E_\pi f^2$.
Independent noninvasive binary readout noise at the two endpoints has the
same correlation. At times $k\Delta$, putting $x=e^{-\lambda\Delta}$ yields
the calibrated moment comparison class already used above. Data processing
of \eqref{eq:ci-diffusion-bridge} proves the same record-KL lower certificate
$\sigma\ge4B/\bar t$.

This corollary composes established likelihood and spectral theory. It is not
a new physical law. It supplies a legitimate mathematical route to suitable
overdamped position measurements, including model classes used for optical
traps; it does not certify that any particular experimental dataset satisfies
the assumptions. Sharpness or finite realization within the CTMC comparison
class does not establish sharpness in a diffusion-only class. The optional
finite-jump activity envelope does not transfer to continuous diffusions.

For an independently computable check, take a two-dimensional rotational
Ornstein--Uhlenbeck process
\[
 dX_t=(-I+\omega J)X_t\,dt+\sqrt2\,dW_t,
 \qquad J=\begin{pmatrix}0&-1\\1&0\end{pmatrix}.
\]
Its stationary covariance is $I$, the reversible drift is $-x$, and
$v=\omega Jx$. Hence $\sigma=2\omega^2$ and the full-path relative entropy
to the reversible reference is $t\omega^2/2$. This unbounded linear drift
does not need a global Novikov assertion: smooth radial cutoffs of the
rotational drift have bounded likelihood controls and preserve the Gaussian
stationary density, since the rotation is orthogonal to radial gradients.
Their expected control energies are uniformly bounded by the Gaussian second
moment. The resulting likelihoods have uniformly bounded relative entropy
and are uniformly integrable. Localization, nonexplosion and uniqueness of
the linear diffusion identify the limiting true likelihood; its finite-energy
Girsanov expression gives the identity above.

Choose the bounded detector mean $f(x)=\sin(x_1)$ and independent binary
readout noise. The stationary mean is zero and
$A=(1-e^{-2})/2$. Since the two scalar Gaussian coordinates at times zero and
$t$ have correlation $\rho=e^{-t}\cos(\omega t)$,
\[
 C(t)=E[f(X_0)f(X_t)]
 =\tfrac12\{e^{-(1-\rho)}-e^{-(1+\rho)}\}
 =e^{-1}\sinh\rho.
\]
The first equality follows directly from the Gaussian characteristic function
and $2\sin x\sin y=\cos(x-y)-\cos(x+y)$.
For $\omega=2$, $t=1$, existing rational exponential and cosine enclosures
certify $C(t)<0$, hence a positive entropy lower certificate, while the same
detector and score yield zero in the $\omega=0$ equilibrium control.
The model entropy rate is exactly eight and the full-path KL is exactly two.
This is an exact Gaussian-model calculation with certified transcendental
enclosures, not a new trajectory experiment or an empirical measurement.

\section{Implementation and checks}
The API \path{constructive_access.calibrated.nonlinear.design} returns the
exact dual lower certificate, a feasible primal upper bracket, and their gap.
The faster \path{certify} accepts supplied rational scores and may reuse them
across independently constructed confidence regions.
\path{confidence_region} constructs the fixed-contrast halfspace and
\path{activity_bridge} supplies the optional activity-calibrated physical step.
All four routines distinguish an assumed observation contract from evidence
that a physical apparatus satisfies it.

The tests compare the rational linear-program solver with an independent
linear optimizer, recover an analytic Bernoulli KL projection, exercise an
equilibrium control, check off-grid normalization and reject malformed primal
points. A positive, decreasing two-lag signal is accepted by the uncalibrated
class but separated after second-moment calibration. Independent scalar root
finding checks the activity inversion, including the equal-bias ring equality
case. Deterministic example outputs are mathematical checks, not measured
laboratory data. Matched held-out performance comparisons are reported in the
integrated benchmark chapter.
